% =====================================================================
% Banco diversificado --- Teoremas de Thevenin e Norton
% 8 N1 + 11 N2 + 11 N3 + 10 N4 = 40 questoes
% Circuitos e valores variados; cada questao pratica possui desenho proprio.
% =====================================================================
\subsubsection*{Banco complementar --- Teoremas de Thévenin e Norton}
\begin{atencao}[Banco visual diversificado]
Os exercícios práticos usam portas, fontes e redes diferentes entre si. A
variação de topologia é intencional: divisores, fontes reais, Norton, escadas,
redes de múltiplas portas, fontes dependentes, caracterização experimental e
cargas não lineares aparecem em problemas distintos.
\end{atencao}

% =====================================================================
% THEVENIN/NORTON --- QUESTOES VISUAIS CRITICAS DOS BANCOS N2 E N3
% Somente desenhos. Enunciados, respostas, resolucoes e niveis permanecem.
% =====================================================================

% N2/8 --- separa o rotulo da carga da anotacao da tensao medida.
\def\fvThNDoisQOitoRevisada{%
\begin{circuitikz}[american,scale=.84,transform shape,line width=.8pt]
  \draw[dashed,cinza,line width=.8pt,rounded corners=6pt]
        (0,0) rectangle (3.4,3.9);
  \node[cinza] at(1.7,2.0){caixa-preta};
  \draw (3.4,3.1)--(5.0,3.1) coordinate(a);
  \draw (3.4,.8)--(5.0,.8) coordinate(b);
  \draw (a) to[R,l_={$R_L=\SI{10}{\ohm}$}] (b);
  \draw[<->,Vcol,line width=.9pt]
        (6.8,1.0)--(6.8,2.9)
        node[midway,right,fill=white,inner sep=1.5pt] {$\SI{24}{\volt}$};
  \node[above,font=\footnotesize,fill=white,inner sep=1.5pt]
        at(1.7,4.15){$V_{oc}=\SI{30}{\volt}$};
\end{circuitikz}}

% N3/9 --- cria um trecho de fio exclusivo para a fonte de 1 A antes do
% resistor de 3 ohms; evita que o ramo vertical termine sobre o resistor.
\def\fvThNTresQNoveRevisada{%
\begin{circuitikz}[american,scale=.58,transform shape,line width=.8pt]
  \coordinate (g) at (6.0,0);
  \coordinate (a) at (6.0,4.8);
  \coordinate (jr) at (10.2,4.8);
  \coordinate (pa) at (11.8,6.3);
  \coordinate (pb) at (11.8,-1.4);
  \draw (0,0) to[\fonteV,l^={$\SI{18}{\volt}$}] (0,4.8)
        to[R,l={$\SI{6}{\ohm}$}] (a);
  \draw (9.0,0) to[isource,l_={$\SI{1}{\ampere}$}] (9.0,4.8);
  \draw (a)--(9.0,4.8)--(jr);
  \draw (13.4,0) to[\fonteV,l_={$\SI{6}{\volt}$}] (13.4,4.8)
        to[R,l_={$\SI{3}{\ohm}$}] (jr);
  \draw (a) to[R,l_={$\SI{6}{\ohm}$}] (g);
  \draw (0,0)--(g)--(9.0,0)--(13.4,0);
  \draw (a)--(6.0,6.3)--(pa) node[ocirc]{} node[right]{A};
  \draw (g)--(6.0,-1.4)--(pb) node[ocirc]{} node[right]{B};
  \fill (a) circle (1.1pt);
  \node[ground] at(g){};
\end{circuitikz}}

% N1 --- Thevenin e Norton: 8 questoes visuais distintas
\blocofixacao

\begin{questao}[1]{}
Determine $V_{Th}$ e $R_{Th}$ vistos em A--B.
\circuitoq{\begin{circuitikz}[american,scale=.86,transform shape]
\coordinate (g) at (3.8,0); \coordinate (a) at (3.8,3.7);
\draw (0,0) to[\fonteV,l^={$\SI{12}{\volt}$}] (0,3.7) to[R,l={$\SI{4}{\ohm}$}] (a);
\draw (a) to[R,l_={$\SI{8}{\ohm}$}] (g)--(0,0);
\draw (a)--(6.2,3.7) node[ocirc]{} node[right]{A}; \draw (g)--(6.2,0) node[ocirc]{} node[right]{B};
\end{circuitikz}}
\resposta{$V_{Th}=\SI{8}{\volt}$; $R_{Th}=\SI{2.67}{\ohm}$}
\resolucao{Em aberto, o divisor fornece $12\cdot8/(4+8)=8$ V. Zerando a fonte, $R_{Th}=4\parallel8=8/3\approx2{,}67\,\Omega$.}
\end{questao}

\begin{questao}[1]{}
Converta o equivalente de Thévenin mostrado em Norton.
\circuitoq{\begin{circuitikz}[american,scale=.9,transform shape]
\draw (0,0) to[\fonteV,l^={$\SI{15}{\volt}$}] (0,3.5) to[R,l={$\SI{5}{\ohm}$}] (4,3.5)--(5.2,3.5) node[ocirc]{};
\draw (0,0)--(5.2,0) node[ocirc]{};
\end{circuitikz}}
\resposta{$I_N=\SI{3}{\ampere}$; $R_N=\SI{5}{\ohm}$}
\resolucao{$I_N=V_{Th}/R_{Th}=15/5=3$ A e $R_N=R_{Th}$.}
\end{questao}

\begin{questao}[1]{}
Determine o equivalente de Thévenin da fonte de Norton.
\circuitoq{\begin{circuitikz}[american,scale=.88,transform shape]
\coordinate (g) at (3.5,0); \coordinate (a) at (3.5,3.6);
\draw (0,0) to[isource,l^={$\SI{3}{\ampere}$}] (0,3.6)--(a);
\draw (a) to[R,l_={$\SI{6}{\ohm}$}] (g)--(0,0);
\draw (a)--(5.5,3.6) node[ocirc]{}; \draw (g)--(5.5,0) node[ocirc]{};
\end{circuitikz}}
\resposta{$V_{Th}=\SI{18}{\volt}$; $R_{Th}=\SI{6}{\ohm}$}
\resolucao{Em aberto toda a corrente passa por $6\,\Omega$, produzindo $18$ V. Com a fonte de corrente aberta, a porta vê apenas $6\,\Omega$.}
\end{questao}

\begin{questao}[1]{}
Para calcular $R_{Th}$, indique o que acontece com as duas fontes independentes.
\circuitoq{\begin{circuitikz}[american,scale=.78,transform shape]
\draw (0,0) to[\fonteV,l^={$E$}] (0,3.7) to[R,l={$R_1$}] (3.5,3.7) coordinate(a);
\draw (a) to[R,l_={$R_2$}] (3.5,0)--(0,0);
\draw (6.8,0) to[isource,l_={$I$}] (6.8,3.7)--(a); \draw (6.8,0)--(3.5,0);
\draw (a)--(8.3,3.7) node[ocirc]{}; \draw (3.5,0)--(8.3,0) node[ocirc]{};
\end{circuitikz}}
\resposta{A fonte de tensão vira curto e a fonte de corrente vira circuito aberto}
\resolucao{Zerar uma fonte ideal significa impor sua grandeza igual a zero: $V=0$ corresponde a curto; $I=0$ corresponde a aberto.}
\end{questao}

\begin{questao}[1]{}
Calcule a corrente e a tensão na carga.
\circuitoq{\begin{circuitikz}[american,scale=.9,transform shape]
\draw (0,0) to[\fonteV,l^={$V_{Th}=\SI{18}{\volt}$}] (0,3.6) to[R,l={$R_{Th}=\SI{3}{\ohm}$}] (4,3.6)
 to[R,l={$R_L=\SI{6}{\ohm}$},i>^={$I_L$}] (4,0)--(0,0);
\end{circuitikz}}
\resposta{$I_L=\SI{2}{\ampere}$; $U_L=\SI{12}{\volt}$}
\resolucao{$I_L=18/(3+6)=2$ A e $U_L=6(2)=12$ V.}
\end{questao}

\begin{questao}[1]{}
Uma rede apresenta $V_{oc}=\SI{20}{\volt}$ e $I_{sc}=\SI{4}{\ampere}$. Determine $R_{Th}$.
\circuitoq{\begin{circuitikz}[american,scale=.85,transform shape]
\draw[dashed,cinza,rounded corners] (0,0) rectangle (3.2,3.5); \node[cinza] at(1.6,1.75){rede};
\draw (3.2,2.8)--(5,2.8) node[ocirc]{}; \draw (3.2,.7)--(5,.7) node[ocirc]{};
\node[right] at(5.3,2.3){$V_{oc}=20$ V}; \node[right] at(5.3,1.3){$I_{sc}=4$ A};
\end{circuitikz}}
\resposta{$R_{Th}=\SI{5}{\ohm}$}
\resolucao{$R_{Th}=V_{oc}/I_{sc}=20/4=5\,\Omega$.}
\end{questao}

\begin{questao}[1]{}
Duas fontes aterradas alimentam a porta. Determine $V_{Th}$ e $R_{Th}$.
\circuitoq{\begin{circuitikz}[american,scale=.74,transform shape]
\coordinate (g) at (5,0); \coordinate (a) at (5,4);
\draw (0,0) to[\fonteV,l^={$\SI{18}{\volt}$}] (0,4) to[R,l={$\SI{6}{\ohm}$}] (a);
\draw (10,0) to[\fonteV,l_={$\SI{6}{\volt}$}] (10,4) to[R,l_={$\SI{3}{\ohm}$}] (a);
\draw (0,0)--(g)--(10,0); \draw (a)--(6.8,4) node[ocirc]{}; \draw (g)--(6.8,0) node[ocirc]{};
\end{circuitikz}}
\resposta{$V_{Th}=\SI{10}{\volt}$; $R_{Th}=\SI{2}{\ohm}$}
\resolucao{Por Millman, $V_{Th}=(18/6+6/3)/(1/6+1/3)=10$ V. Zerando as fontes, $R_{Th}=6\parallel3=2\,\Omega$.}
\end{questao}

\begin{questao}[1]{}
Determine o equivalente visto na porta.
\circuitoq{\begin{circuitikz}[american,scale=.84,transform shape]
\coordinate (g) at (4,0); \coordinate (a) at (4,3.8);
\draw (0,0) to[\fonteV,l^={$\SI{30}{\volt}$}] (0,3.8) to[R,l={$\SI{10}{\ohm}$}] (a);
\draw (a) to[R,l_={$\SI{15}{\ohm}$}] (g)--(0,0);
\draw (a)--(6.2,3.8) node[ocirc]{}; \draw (g)--(6.2,0) node[ocirc]{};
\end{circuitikz}}
\resposta{$V_{Th}=\SI{18}{\volt}$; $R_{Th}=\SI{6}{\ohm}$}
\resolucao{$V_{Th}=30\cdot15/(10+15)=18$ V e $R_{Th}=10\parallel15=6\,\Omega$.}
\end{questao}


% N2: substitui somente a oitava figura, na qual a anotacao de 24 V colidia
% com o rotulo da carga.
\begingroup
\newcount\fvThNDoisCircuitCount
\fvThNDoisCircuitCount=0
\let\fvThNDoisCircuitoOriginal\circuitoq
\renewcommand{\circuitoq}[1]{%
  \global\advance\fvThNDoisCircuitCount by 1\relax
  \ifnum\fvThNDoisCircuitCount=8\relax
    \fvThNDoisCircuitoOriginal{\fvThNDoisQOitoRevisada}%
  \else
    \fvThNDoisCircuitoOriginal{#1}%
  \fi}
% N2 --- Thevenin e Norton: 11 questoes com topologias distintas
% Revisao visual: portas roteadas fora dos ramos e maior separacao de rotulos.
\blocoaplicacao

\begin{questao}[2]{}
Determine o equivalente de Thévenin visto de A--B.
\circuitoq{\begin{circuitikz}[american,scale=.86,transform shape,line width=.8pt]
\coordinate (g) at (4.2,0); \coordinate (a) at (4.2,4.0);
\draw (0,0) to[\fonteV,l^={$\SI{24}{\volt}$}] (0,4.0) to[R,l={$\SI{6}{\kilo\ohm}$}] (a);
\draw (a) to[R,l_={$\SI{3}{\kilo\ohm}$}] (g)--(0,0);
\draw (a)--(6.6,4.0) node[ocirc]{} node[right]{A};
\draw (g)--(6.6,0) node[ocirc]{} node[right]{B};
\fill (a) circle (1.1pt); \node[ground] at(g){};
\end{circuitikz}}
\resposta{$V_{Th}=\SI{8}{\volt}$; $R_{Th}=\SI{2}{\kilo\ohm}$}
\resolucao{Divisor em vazio: $24\cdot3/(6+3)=8$ V. Com a fonte em curto, $6\parallel3=2$ k$\Omega$.}
\end{questao}

\begin{questao}[2]{}
Determine $V_{Th}$ e $R_{Th}$ vistos na porta à direita.
\circuitoq{\begin{circuitikz}[american,scale=.82,transform shape,line width=.8pt]
\coordinate (g) at (2.8,0); \coordinate (a) at (2.8,4.0); \coordinate (b) at (7.6,4.0);
\draw (0,0) to[isource,l^={$\SI{2}{\milli\ampere}$}] (0,4.0)--(a);
\draw (a) to[R,l_={$\SI{4}{\kilo\ohm}$}] (g)--(0,0);
\draw (a) to[R,l={$\SI{2}{\kilo\ohm}$}] (b) node[ocirc]{} node[right]{A};
\draw (g)--(7.6,0) node[ocirc]{} node[right]{B};
\fill (a) circle (1.1pt); \node[ground] at(g){};
\end{circuitikz}}
\resposta{$V_{Th}=\SI{8}{\volt}$; $R_{Th}=\SI{6}{\kilo\ohm}$}
\resolucao{Em aberto não circula corrente no resistor de 2 k$\Omega$, então $V_{Th}=2$ mA$\times4$ k$\Omega=8$ V. Com a fonte aberta, a porta vê $2+4=6$ k$\Omega$.}
\end{questao}

\begin{questao}[2]{}
Duas fontes reais aterradas alimentam a mesma porta. Determine o equivalente.
\circuitoq{\begin{circuitikz}[american,scale=.68,transform shape,line width=.8pt]
\coordinate (g) at (5.4,0); \coordinate (a) at (5.4,4.4);
\coordinate (pa) at (7.2,5.8); \coordinate (pb) at (7.2,-1.3);
\draw (0,0) to[\fonteV,l^={$\SI{20}{\volt}$}] (0,4.4) to[R,l={$\SI{5}{\kilo\ohm}$}] (a);
\draw (11.0,0) to[\fonteV,l_={$\SI{10}{\volt}$}] (11.0,4.4) to[R,l_={$\SI{10}{\kilo\ohm}$}] (a);
\draw (0,0)--(g)--(11.0,0);
\draw (a)--(5.4,5.8)--(pa) node[ocirc]{} node[right]{A};
\draw (g)--(5.4,-1.3)--(pb) node[ocirc]{} node[right]{B};
\fill (a) circle (1.1pt); \node[ground] at(g){};
\end{circuitikz}}
\resposta{$V_{Th}\approx\SI{16.67}{\volt}$; $R_{Th}\approx\SI{3.33}{\kilo\ohm}$}
\resolucao{$V_{Th}=(20/5+10/10)/(1/5+1/10)=16{,}67$ V. Zerando as fontes, $R_{Th}=5\parallel10=3{,}33$ k$\Omega$.}
\end{questao}

\begin{questao}[2]{}
A porta está no nó B. Determine seu equivalente de Thévenin.
\circuitoq{\begin{circuitikz}[american,scale=.74,transform shape,line width=.8pt]
\coordinate (g) at (3.8,0); \coordinate (a) at (3.8,4.2); \coordinate (b) at (9.2,4.2);
\draw (0,0) to[\fonteV,l^={$\SI{12}{\volt}$}] (0,4.2) to[R,l={$R_1=\SI{2}{\kilo\ohm}$}] (a);
\draw (a) to[R,l_={$R_2=\SI{4}{\kilo\ohm}$}] (g)--(0,0);
\draw (a) to[R,l={$R_3=\SI{3}{\kilo\ohm}$}] (b) node[ocirc]{} node[right]{B};
\draw (g)--(9.2,0) node[ocirc]{} node[right]{terra};
\fill (a) circle (1.1pt); \node[Vcol,fill=white,inner sep=1.4pt] at ($(a)+(0,.65)$) {$A$};
\node[ground] at(g){};
\end{circuitikz}}
\resposta{$V_{Th}=\SI{8}{\volt}$; $R_{Th}\approx\SI{4.33}{\kilo\ohm}$}
\resolucao{Em aberto, $R_3$ não conduz e $V_B=12\cdot4/(2+4)=8$ V. Com a fonte zerada, $R_{Th}=R_3+(R_1\parallel R_2)=3+4/3=4{,}33$ k$\Omega$.}
\end{questao}

\begin{questao}[2]{}
Determine o equivalente de Thévenin da fonte de Norton com resistor série até a porta.
\circuitoq{\begin{circuitikz}[american,scale=.82,transform shape,line width=.8pt]
\coordinate (g) at (3.0,0); \coordinate (a) at (3.0,4.0); \coordinate (b) at (7.8,4.0);
\draw (0,0) to[isource,l^={$\SI{5}{\milli\ampere}$}] (0,4.0)--(a);
\draw (a) to[R,l_={$\SI{3}{\kilo\ohm}$}] (g)--(0,0);
\draw (a) to[R,l={$\SI{2}{\kilo\ohm}$}] (b) node[ocirc]{} node[right]{A};
\draw (g)--(7.8,0) node[ocirc]{} node[right]{B};
\fill (a) circle (1.1pt); \node[ground] at(g){};
\end{circuitikz}}
\resposta{$V_{Th}=\SI{15}{\volt}$; $R_{Th}=\SI{5}{\kilo\ohm}$}
\resolucao{Em vazio o resistor série de 2 k$\Omega$ não conduz, então $V_{Th}=5$ mA$\times3$ k$\Omega=15$ V. Com a fonte aberta, $R_{Th}=3+2=5$ k$\Omega$.}
\end{questao}

\begin{questao}[2]{}
Determine $V_{Th}$ e $R_{Th}$ vistos em A--B.
\circuitoq{\begin{circuitikz}[american,scale=.70,transform shape,line width=.8pt]
\coordinate (g) at (4.8,0); \coordinate (a) at (4.8,4.4);
\coordinate (pa) at (9.8,5.7); \coordinate (pb) at (9.8,-1.3);
\draw (0,0) to[\fonteV,l^={$\SI{12}{\volt}$}] (0,4.4) to[R,l={$\SI{4}{\ohm}$}] (a);
\draw (a) to[R,l_={$\SI{6}{\ohm}$}] (g)--(0,0);
\draw (8.4,0) to[isource,l_={$\SI{2}{\ampere}$}] (8.4,4.4)--(a); \draw (8.4,0)--(g);
\draw (a)--(4.8,5.7)--(pa) node[ocirc]{} node[right]{A};
\draw (g)--(4.8,-1.3)--(pb) node[ocirc]{} node[right]{B};
\fill (a) circle (1.1pt); \node[ground] at(g){};
\end{circuitikz}}
\resposta{$V_{Th}=\SI{12}{\volt}$; $R_{Th}=\SI{2.4}{\ohm}$}
\resolucao{A condutância da porta é $1/4+1/6=5/12$ S. A fonte de 12 V equivale a 3 A e soma-se à fonte de 2 A: $V_{Th}=5/(5/12)=12$ V. Zerando as fontes, $R_{Th}=4\parallel6=2{,}4\,\Omega$.}
\end{questao}

\begin{questao}[2]{}
O equivalente abaixo alimentará três cargas distintas. Complete as correntes e identifique a condição de maior potência.
\circuitoq{\begin{circuitikz}[american,scale=.86,transform shape,line width=.8pt]
\draw (0,0) to[\fonteV,l^={$V_{Th}=\SI{9}{\volt}$}] (0,3.8)
      to[R,l={$R_{Th}=\SI{3}{\ohm}$}] (4.2,3.8)
      to[R,l={$R_L$}] (4.2,0)--(0,0);
\node[right,align=left,font=\footnotesize,draw=cinza,rounded corners,inner sep=3pt]
      at(6.3,1.9){$R_L=3\,\Omega$\\$R_L=6\,\Omega$\\$R_L=15\,\Omega$};
\end{circuitikz}}
\resposta{$I_L=\SI{1.5}{\ampere}$, $\SI{1}{\ampere}$ e $\SI{0.5}{\ampere}$; maior potência em $R_L=\SI{3}{\ohm}$}
\resolucao{$I_L=9/(3+R_L)$. As correntes são 1,5 A, 1 A e 0,5 A. A máxima transferência ocorre quando $R_L=R_{Th}=3\,\Omega$.}
\end{questao}

\begin{questao}[2]{}
Uma caixa-preta fornece $\SI{30}{\volt}$ em vazio e $\SI{24}{\volt}$ com $R_L=\SI{10}{\ohm}$. Determine o equivalente e $I_{sc}$.
\circuitoq{\begin{circuitikz}[american,scale=.84,transform shape,line width=.8pt]
\draw[dashed,cinza,line width=.8pt,rounded corners=6pt] (0,0) rectangle (3.4,3.9);
\node[cinza] at(1.7,2.0){caixa-preta};
\draw (3.4,3.1)--(5.0,3.1) coordinate(a); \draw (3.4,.8)--(5.0,.8) coordinate(b);
\draw (a) to[R,l={$R_L=\SI{10}{\ohm}$}] (b);
\draw[<->,Vcol,line width=.9pt] (6.0,1.0)--(6.0,2.9) node[midway,right] {$\SI{24}{\volt}$};
\node[above,font=\footnotesize,fill=white,inner sep=1.5pt] at(1.7,4.15){$V_{oc}=\SI{30}{\volt}$};
\end{circuitikz}}
\resposta{$V_{Th}=\SI{30}{\volt}$; $R_{Th}=\SI{2.5}{\ohm}$; $I_{sc}=\SI{12}{\ampere}$}
\resolucao{Com a carga, $I=24/10=2{,}4$ A e a queda interna é 6 V; logo $R_{Th}=6/2{,}4=2{,}5\,\Omega$. Então $I_{sc}=30/2{,}5=12$ A.}
\end{questao}

\begin{questao}[2]{}
A mesma rede é observada em duas portas. Determine os equivalentes em A e B.
\circuitoq{\begin{circuitikz}[american,scale=.68,transform shape,line width=.8pt]
\coordinate (g) at (4.0,0); \coordinate (a) at (4.0,4.5); \coordinate (b) at (10.0,4.5);
\coordinate (pa) at (6.0,6.0);
\draw (0,0) to[\fonteV,l^={$\SI{18}{\volt}$}] (0,4.5) to[R,l={$\SI{6}{\kilo\ohm}$}] (a);
\draw (a) to[R,l_={$\SI{3}{\kilo\ohm}$}] (g)--(0,0);
\draw (a) to[R,l={$\SI{4}{\kilo\ohm}$}] (b) node[ocirc]{} node[right]{B};
\draw (a)--(4.0,6.0)--(pa) node[ocirc]{} node[right]{A};
\draw (g)--(10.0,0) node[ocirc]{} node[right]{terra};
\fill (a) circle (1.1pt); \node[ground] at(g){};
\end{circuitikz}}
\resposta{Porta A: $V_{Th}=\SI{6}{\volt}$, $R_{Th}=\SI{2}{\kilo\ohm}$; porta B: $V_{Th}=\SI{6}{\volt}$, $R_{Th}=\SI{6}{\kilo\ohm}$}
\resolucao{Em vazio, nenhuma corrente passa no resistor até B, então ambas as portas têm 6 V. Em A, $R_{Th}=6\parallel3=2$ k$\Omega$; em B soma-se o resistor de 4 k$\Omega$, resultando 6 k$\Omega$.}
\end{questao}

\begin{questao}[2]{}
Determine o equivalente de Norton da rede.
\circuitoq{\begin{circuitikz}[american,scale=.84,transform shape,line width=.8pt]
\coordinate (g) at (4.2,0); \coordinate (a) at (4.2,4.0);
\draw (0,0) to[isource,l^={$\SI{4}{\ampere}$}] (0,4.0)--(a);
\draw (a) to[R,l_={$\SI{12}{\ohm}$}] (2.8,0)--(g);
\draw (a) to[R,l={$\SI{6}{\ohm}$}] (5.8,0)--(g); \draw (0,0)--(g);
\draw (a)--(7.8,4.0) node[ocirc]{} node[right]{A};
\draw (g)--(7.8,0) node[ocirc]{} node[right]{B};
\node[ground] at(g){};
\end{circuitikz}}
\resposta{$I_N=\SI{4}{\ampere}$; $R_N=\SI{4}{\ohm}$}
\resolucao{Em curto, toda a corrente da fonte vai para a porta, logo $I_N=4$ A. Com a fonte aberta, $R_N=12\parallel6=4\,\Omega$.}
\end{questao}

\begin{questao}[2]{}
Uma carga $R_L=\SI{5}{\ohm}$ é conectada ao equivalente de Norton mostrado. Determine $I_L$ e $U_L$.
\circuitoq{\begin{circuitikz}[american,scale=.86,transform shape,line width=.8pt]
\coordinate (g) at (4.2,0); \coordinate (a) at (4.2,4.0);
\draw (0,0) to[isource,l^={$I_N=\SI{3}{\ampere}$}] (0,4.0)--(a);
\draw (a) to[R,l_={$R_N=\SI{6}{\ohm}$}] (2.8,0)--(g);
\draw (a) to[R,l={$R_L=\SI{5}{\ohm}$},i>^={$I_L$}] (6.2,0)--(g);
\draw (0,0)--(g); \node[ground] at(g){};
\end{circuitikz}}
\resposta{$I_L\approx\SI{1.64}{\ampere}$; $U_L\approx\SI{8.18}{\volt}$}
\resolucao{Por divisão de corrente, $I_L=3\cdot6/(6+5)=18/11\approx1{,}64$ A. Então $U_L=5I_L\approx8{,}18$ V.}
\end{questao}
\endgroup

% N3: substitui somente a nona figura, separando o ramo da fonte de corrente
% do resistor ligado a fonte de tensao da direita.
\begingroup
\newcount\fvThNTresCircuitCount
\fvThNTresCircuitCount=0
\let\fvThNTresCircuitoOriginal\circuitoq
\renewcommand{\circuitoq}[1]{%
  \global\advance\fvThNTresCircuitCount by 1\relax
  \ifnum\fvThNTresCircuitCount=9\relax
    \fvThNTresCircuitoOriginal{\fvThNTresQNoveRevisada}%
  \else
    \fvThNTresCircuitoOriginal{#1}%
  \fi}
% N3 --- Thevenin e Norton: 11 questoes avancadas e distintas
\blocoaprofundamento

\begin{questao}[3]{}
A rede contém apenas uma fonte dependente. Determine $R_{Th}$ pela fonte de teste.
\circuitoq{\begin{circuitikz}[american,scale=.8,transform shape]
\coordinate (g) at (3.6,0); \coordinate (a) at (3.6,4);
\draw (a) to[R,l_={$\SI{1}{\kilo\ohm}$}] (g);
\draw (0,0) to[cisource,l^={$0{,}5\,\mathrm{mS}\,V_A$}] (0,4)--(a); \draw (0,0)--(g);
\draw (7.3,0) to[isource,l_={$I_t$}] (7.3,4)--(a); \draw (7.3,0)--(g);
\node[above,Vcol] at(a){$V_A$}; \node[ground] at(g){};
\end{circuitikz}}
\resposta{$R_{Th}=\SI{2}{\kilo\ohm}$}
\resolucao{$I_t=V_A/1000-0{,}0005V_A=0{,}0005V_A$. Logo $R_{Th}=V_A/I_t=2$ k$\Omega$.}
\end{questao}

\begin{questao}[3]{}
Determine $V_{Th}$, $I_{sc}$ e $R_{Th}$, mantendo a fonte dependente ativa.
\circuitoq{\begin{circuitikz}[american,scale=.72,transform shape]
\coordinate (g) at (4.2,0); \coordinate (a) at (4.2,4.1);
\draw (0,0) to[\fonteV,l^={$\SI{12}{\volt}$}] (0,4.1) to[R,l={$\SI{2}{\kilo\ohm}$}] (a);
\draw (a) to[R,l_={$\SI{2}{\kilo\ohm}$}] (g)--(0,0);
\draw (8,0) to[cisource,l_={$0{,}25\,\mathrm{mS}\,V_A$}] (8,4.1)--(a); \draw (8,0)--(g);
\draw (a)--(10,4.1) node[ocirc]{}; \draw (g)--(10,0) node[ocirc]{};
\node[above,Vcol] at(a){$V_A$}; \node[ground] at(g){};
\end{circuitikz}}
\resposta{$V_{Th}=\SI{8}{\volt}$; $I_{sc}=\SI{6}{\milli\ampere}$; $R_{Th}\approx\SI{1.33}{\kilo\ohm}$}
\resolucao{Em aberto, $(V-12)/2k+V/2k-0{,}25\,\mathrm{mS}\,V=0$, dando $V=8$ V. Em curto, $V=0$ e a corrente é $12/2k=6$ mA. Logo $R_{Th}=8/6$ k$\Omega\approx1{,}33$ k$\Omega$.}
\end{questao}

\begin{questao}[3]{}
Uma rede linear foi medida com duas cargas. Determine $V_{Th}$ e $R_{Th}$ sem usar circuito aberto.
\circuitoq{\begin{circuitikz}[american,scale=.82,transform shape]
\draw[dashed,cinza,rounded corners] (0,0) rectangle (3.2,3.8); \node[cinza] at(1.6,2){rede};
\draw (3.2,3.1)--(4.8,3.1) coordinate(a); \draw (3.2,.7)--(4.8,.7) coordinate(b);
\draw (a) to[R,l={$R_L$}] (b);
\node[right,align=left,font=\footnotesize] at(5.7,2){$R_L=2\,k\Omega\Rightarrow U=6$ V\\$R_L=6\,k\Omega\Rightarrow U=9$ V};
\end{circuitikz}}
\resposta{$V_{Th}=\SI{12}{\volt}$; $R_{Th}=\SI{2}{\kilo\ohm}$}
\resolucao{Usando $U=V_{Th}R_L/(R_{Th}+R_L)$ nas duas medições, obtém-se o sistema que resulta em $R_{Th}=2$ k$\Omega$ e $V_{Th}=12$ V.}
\end{questao}

\begin{questao}[3]{}
No divisor, ambos os resistores têm tolerância de $\pm5\%$. Determine os extremos de $V_{Th}$ e de $R_{Th}$.
\circuitoq{\begin{circuitikz}[american,scale=.82,transform shape]
\coordinate (g) at (4,0); \coordinate (a) at (4,3.8);
\draw (0,0) to[\fonteV,l^={$\SI{12}{\volt}$}] (0,3.8) to[R,l={$6\,k\Omega\;\pm5\%$}] (a);
\draw (a) to[R,l_={$3\,k\Omega\;\pm5\%$}] (g)--(0,0); \draw (a)--(6.2,3.8) node[ocirc]{}; \draw (g)--(6.2,0) node[ocirc]{};
\end{circuitikz}}
\resposta{$V_{Th}\approx\SI{3.74}{\volt}$ a $\SI{4.27}{\volt}$; $R_{Th}=\SI{1.90}{\kilo\ohm}$ a $\SI{2.10}{\kilo\ohm}$}
\resolucao{$V_{Th}$ máximo ocorre com $R_2$ alto e $R_1$ baixo; o mínimo no caso oposto. Já $R_{Th}=R_1\parallel R_2$ varia de 1,90 a 2,10 k$\Omega$.}
\end{questao}

\begin{questao}[3]{}
A fonte não pode ser curto-circuitada. A partir de duas medições, estime seu equivalente.
\circuitoq{\begin{circuitikz}[american,scale=.82,transform shape]
\draw[dashed,cinza,rounded corners] (0,0) rectangle (3.3,3.8); \node[cinza] at(1.65,2){fonte};
\draw (3.3,3.1)--(4.8,3.1) coordinate(a); \draw (3.3,.7)--(4.8,.7) coordinate(b); \draw (a) to[R,l={$R_L$},i>^={$I$}] (b);
\node[right,align=left,font=\footnotesize] at(5.8,2){$5{,}5\,\Omega\Rightarrow10$ A\\$2{,}5\,\Omega\Rightarrow20$ A};
\end{circuitikz}}
\resposta{$V_{Th}=\SI{60}{\volt}$; $R_{Th}=\SI{0.5}{\ohm}$; $I_N=\SI{120}{\ampere}$}
\resolucao{As tensões são 55 V e 50 V. Pela reta $U=V_{Th}-R_{Th}I$, a inclinação é $-(50-55)/(20-10)=-0{,}5$, então $R_{Th}=0{,}5\,\Omega$ e $V_{Th}=60$ V.}
\end{questao}

\begin{questao}[3]{}
O equivalente alimenta um LED aproximado por queda constante de $\SI{2}{\volt}$. Determine o ponto de operação.
\circuitoq{\begin{circuitikz}[american,scale=.88,transform shape]
\draw (0,0) to[\fonteV,l^={$V_{Th}=\SI{5}{\volt}$}] (0,3.7) to[R,l={$R_{Th}=\SI{200}{\ohm}$}] (4,3.7)
 to[leD,l={$LED\;(2\,V)$},i>^={$I$}] (4,0)--(0,0);
\end{circuitikz}}
\resposta{$I=\SI{15}{\milli\ampere}$}
\resolucao{A reta de carga é $U=5-200I$. Na aproximação $U_D=2$ V: $I=(5-2)/200=15$ mA.}
\end{questao}

\begin{questao}[3]{}
A carga deve receber no mínimo $\SI{20}{\volt}$. Determine a faixa admissível de $R_L$.
\circuitoq{\begin{circuitikz}[american,scale=.88,transform shape]
\draw (0,0) to[\fonteV,l^={$V_{Th}=\SI{24}{\volt}$}] (0,3.7) to[R,l={$R_{Th}=\SI{8}{\ohm}$}] (4,3.7)
 to[R,l={$R_L$}] (4,0)--(0,0); \draw[<->,Vcol] (4.9,.4)--(4.9,3.3) node[midway,right] {$U_L\ge20$ V};
\end{circuitikz}}
\resposta{$R_L\ge\SI{40}{\ohm}$}
\resolucao{$24R_L/(R_L+8)\ge20$ implica $24R_L\ge20R_L+160$, portanto $R_L\ge40\,\Omega$.}
\end{questao}

\begin{questao}[3]{}
Compare as duas portas da mesma rede para uma carga de $\SI{4}{\kilo\ohm}$.
\circuitoq{\begin{circuitikz}[american,scale=.72,transform shape]
\coordinate (g) at (3.4,0); \coordinate (a) at (3.4,4); \coordinate (b) at (8.5,4);
\draw (0,0) to[\fonteV,l^={$\SI{18}{\volt}$}] (0,4) to[R,l={$\SI{6}{\kilo\ohm}$}] (a);
\draw (a) to[R,l_={$\SI{3}{\kilo\ohm}$}] (g)--(0,0); \draw (a)--(5,4) node[ocirc]{} node[above]{A};
\draw (a) to[R,l={$\SI{4}{\kilo\ohm}$}] (b) node[ocirc]{} node[right]{B}; \draw (g)--(8.5,0) node[ocirc]{};
\end{circuitikz}}
\resposta{Porta A: $U_L=\SI{4}{\volt}$; porta B: $U_L=\SI{2.4}{\volt}$}
\resolucao{Em A, o equivalente é 6 V em série com 2 k$\Omega$: $U=6\cdot4/(2+4)=4$ V. Em B, é 6 V em série com 6 k$\Omega$: $U=6\cdot4/(6+4)=2{,}4$ V.}
\end{questao}

\begin{questao}[3]{}
Determine o equivalente visto pela porta da rede com três fontes independentes.
\circuitoq{\begin{circuitikz}[american,scale=.68,transform shape]
\coordinate (g) at (5,0); \coordinate (a) at (5,4.2);
\draw (0,0) to[\fonteV,l^={$\SI{18}{\volt}$}] (0,4.2) to[R,l={$\SI{6}{\ohm}$}] (a);
\draw (10,0) to[\fonteV,l_={$\SI{6}{\volt}$}] (10,4.2) to[R,l_={$\SI{3}{\ohm}$}] (a);
\draw (7.8,0) to[isource,l_={$\SI{1}{\ampere}$}] (7.8,4.2)--(a); \draw (a) to[R,l_={$\SI{6}{\ohm}$}] (g);
\draw (0,0)--(g)--(7.8,0)--(10,0); \draw (a)--(11.5,4.2) node[ocirc]{}; \draw (g)--(11.5,0) node[ocirc]{};
\end{circuitikz}}
\resposta{$V_{Th}=\SI{9}{\volt}$; $R_{Th}=\SI{1.5}{\ohm}$}
\resolucao{A soma Norton equivalente é $18/6+6/3+1=6$ A e a condutância é $1/6+1/3+1/6=2/3$ S. Logo $V_{Th}=9$ V e $R_{Th}=1/(2/3)=1{,}5\,\Omega$.}
\end{questao}

\begin{questao}[3]{}
Três medições fornecem os pares $(I,U)=(1,10)$, $(2,8)$ e $(3,6)$ em unidades SI. Determine o equivalente e avalie a linearidade.
\circuitoq{\begin{circuitikz}[american,scale=.82,transform shape]
\draw[dashed,cinza,rounded corners] (0,0) rectangle (3.3,3.7); \node[cinza] at(1.65,1.85){rede};
\draw (3.3,3)--(4.8,3) node[ocirc]{}; \draw (3.3,.7)--(4.8,.7) node[ocirc]{};
\node[right,align=left,font=\footnotesize] at(5.3,1.9){$(1\,A,10\,V)$\\$(2\,A,8\,V)$\\$(3\,A,6\,V)$};
\end{circuitikz}}
\resposta{$V_{Th}=\SI{12}{\volt}$; $R_{Th}=\SI{2}{\ohm}$; dados perfeitamente lineares}
\resolucao{Os três pontos pertencem à reta $U=12-2I$. O intercepto é $V_{Th}=12$ V e a inclinação é $-R_{Th}=-2\,\Omega$.}
\end{questao}

\begin{questao}[3]{}
A rede ativa abaixo pode apresentar resistência de Thévenin negativa. Determine-a com uma fonte de teste.
\circuitoq{\begin{circuitikz}[american,scale=.8,transform shape]
\coordinate (g) at (3.6,0); \coordinate (a) at (3.6,4);
\draw (a) to[R,l_={$\SI{1}{\kilo\ohm}$}] (g);
\draw (0,0) to[cisource,l^={$1{,}5\,\mathrm{mS}\,V_A$}] (0,4)--(a); \draw (0,0)--(g);
\draw (7.2,0) to[isource,l_={$I_t$}] (7.2,4)--(a); \draw (7.2,0)--(g); \node[ground] at(g){};
\end{circuitikz}}
\resposta{$R_{Th}=\SI{-2}{\kilo\ohm}$}
\resolucao{$I_t=V/1000-1{,}5\,\mathrm{mS}\,V=-0{,}5\,\mathrm{mS}\,V$, logo $V/I_t=-2$ k$\Omega$. A rede ativa fornece energia à porta.}
\end{questao}

\endgroup

% N4 --- Thevenin e Norton: 10 desafios
\blocodesafio

\begin{questao}[4]{}
Prove o teorema de Thévenin aplicando uma fonte de corrente de teste $I$ à porta e usando superposição.
\circuitoq{\begin{circuitikz}[american,scale=.82,transform shape]
\draw[dashed,cinza,rounded corners] (0,0) rectangle (3.4,3.8); \node[cinza] at(1.7,1.9){rede linear};
\draw (3.4,3.1)--(5,3.1) coordinate(a); \draw (3.4,.7)--(5,.7) coordinate(b);
\draw (b) to[isource,l_={$I$}] (a); \draw[<->,Vcol] (5.7,.9)--(5.7,2.9) node[midway,right] {$U$};
\end{circuitikz}}
\resposta{$U=V_{oc}-R_{Th}I$}
\resolucao{Por superposição, $U=U'+U''$. Com a fonte de teste zerada, a porta está aberta e $U'=V_{oc}$. Com as fontes internas zeradas, a rede vista pela porta é linear e produz $U''=-R_{Th}I$.}
\end{questao}

\begin{questao}[4]{}
Mostre que os modelos abaixo são equivalentes para qualquer carga e obtenha as relações entre seus parâmetros.
\circuitoq{\begin{circuitikz}[american,scale=.7,transform shape]
\draw (0,0) to[\fonteV,l^={$V_{Th}$}] (0,3.4) to[R,l={$R_{Th}$}] (3.5,3.4)--(4.5,3.4) node[ocirc]{};
\draw (0,0)--(4.5,0) node[ocirc]{}; \node[above] at(2.2,4){Thévenin};
\draw (7,0) to[isource,l^={$I_N$}] (7,3.4)--(10.3,3.4)--(11.3,3.4) node[ocirc]{};
\draw (10.3,3.4) to[R,l_={$R_N$}] (10.3,0); \draw (7,0)--(11.3,0) node[ocirc]{}; \node[above] at(9.1,4){Norton};
\end{circuitikz}}
\resposta{$R_N=R_{Th}$ e $I_N=V_{Th}/R_{Th}$}
\resolucao{Thévenin impõe $U=V_{Th}-R_{Th}I$. Norton impõe $U=I_NR_N-R_NI$. As duas retas coincidem quando intercepto e inclinação são iguais.}
\end{questao}

\begin{questao}[4]{}
Deduza um método de mínimos quadrados para estimar $V_{Th}$ e $R_{Th}$ a partir de $n\ge3$ medições $(I_k,U_k)$.
\resposta{Ajustar a reta $U=V_{Th}-R_{Th}I$; intercepto $=V_{Th}$ e inclinação $=-R_{Th}$}
\resolucao{Use regressão linear: $-R_{Th}=\sum(I_k-\bar I)(U_k-\bar U)/\sum(I_k-\bar I)^2$ e $V_{Th}=\bar U+R_{Th}\bar I$. Resíduos sistemáticos indicam não linearidade.}
\end{questao}

\begin{questao}[4]{}
Projete um ensaio seguro para estimar o Norton de uma fonte de $\SI{60}{\volt}$ cuja corrente de curto prevista é $\SI{120}{\ampere}$, usando instrumentos limitados a $\SI{25}{\ampere}$.
\circuitoq{\begin{circuitikz}[american,scale=.82,transform shape]
\draw[dashed,cinza,rounded corners] (0,0) rectangle (3.3,3.8); \node[cinza] at(1.65,2){fonte};
\draw (3.3,3.1)--(4.8,3.1) coordinate(a); \draw (3.3,.7)--(4.8,.7) coordinate(b); \draw (a) to[R,l={$R_L$},i>^={$I$}] (b);
\node[right,align=left,font=\footnotesize] at(5.7,2){usar dois ou mais\\pontos abaixo de 25 A};
\end{circuitikz}}
\resposta{Exemplo: $R_L=\SI{5.5}{\ohm}$ e $\SI{2.5}{\ohm}$, produzindo aproximadamente $10$ A e $20$ A para $R_{Th}\approx0{,}5\,\Omega$}
\resolucao{Escolhem-se cargas que mantenham a corrente abaixo do limite e suficientemente separadas para evitar mau condicionamento. Ajusta-se a reta $U=V_{Th}-R_{Th}I$ e extrapola-se $I_N=V_{Th}/R_{Th}$ sem curto real.}
\end{questao}

\begin{questao}[4]{}
A rede de duas portas não pode ser representada simultaneamente por dois equivalentes independentes se houver acoplamento. Explique a necessidade de uma matriz de porta.
\circuitoq{\begin{circuitikz}[american,scale=.82,transform shape]
\draw[dashed,cinza,rounded corners] (0,0) rectangle (4.2,4); \node[cinza] at(2.1,2){rede acoplada};
\draw (0,3.2)--(-1.4,3.2) node[ocirc]{} node[left]{A+}; \draw (0,.8)--(-1.4,.8) node[ocirc]{} node[left]{A-};
\draw (4.2,3.2)--(5.6,3.2) node[ocirc]{} node[right]{B+}; \draw (4.2,.8)--(5.6,.8) node[ocirc]{} node[right]{B-};
\end{circuitikz}}
\resposta{É necessário um modelo de dois acessos, por exemplo $\mathbf U=\mathbf V_0-\mathbf R\mathbf I$}
\resolucao{Uma corrente aplicada em A pode alterar a tensão em B e vice-versa. Os termos fora da diagonal da matriz representam esse acoplamento, algo que dois Thévenins independentes não capturam.}
\end{questao}

\begin{questao}[4]{}
Para o equivalente abaixo, maximize a potência na carga sob a restrição $U_L\ge\SI{20}{\volt}$.
\circuitoq{\begin{circuitikz}[american,scale=.88,transform shape]
\draw (0,0) to[\fonteV,l^={$\SI{24}{\volt}$}] (0,3.8) to[R,l={$\SI{8}{\ohm}$}] (4.2,3.8) to[R,l={$R_L$}] (4.2,0)--(0,0);
\draw[<->,Vcol] (5,.4)--(5,3.4) node[midway,right] {$U_L\ge20$ V};
\end{circuitikz}}
\resposta{$R_L=\SI{40}{\ohm}$ no limite; $P_L=\SI{10}{\watt}$}
\resolucao{A restrição exige $R_L\ge40\,\Omega$. Para $R_L>R_{Th}$, a potência decresce com $R_L$, portanto o máximo admissível ocorre no menor valor permitido, 40 $\Omega$.}
\end{questao}

\begin{questao}[4]{}
Aplique uma fonte de teste e determine quando a rede ativa abaixo apresenta $R_{Th}<0$.
\circuitoq{\begin{circuitikz}[american,scale=.78,transform shape]
\coordinate (g) at (3.6,0); \coordinate (a) at (3.6,4);
\draw (a) to[R,l_={$R$}] (g); \draw (0,0) to[cisource,l^={$g_mV_A$}] (0,4)--(a); \draw (0,0)--(g);
\draw (7.2,0) to[isource,l_={$I_t$}] (7.2,4)--(a); \draw (7.2,0)--(g); \node[ground] at(g){};
\end{circuitikz}}
\resposta{$R_{Th}=1/(1/R-g_m)$; é negativo quando $g_m>1/R$}
\resolucao{$I_t=V_A/R-g_mV_A$. Logo $R_{Th}=V_A/I_t=1/(1/R-g_m)$. Se a transcondutância superar a condutância passiva, a inclinação da porta muda de sinal.}
\end{questao}

\begin{questao}[4]{}
Dois divisores entregam a mesma tensão em vazio, mas têm resistências de Thévenin diferentes. Formule um critério de projeto para escolher a porta mais robusta à variação de carga.
\resposta{Preferir a menor $R_{Th}$ para manter a tensão mais próxima de $V_{Th}$ sob carga}
\resolucao{$U_L=V_{Th}R_L/(R_{Th}+R_L)$. Para um mesmo $V_{Th}$ e $R_L$, a queda relativa cresce monotonicamente com $R_{Th}$. Portanto a porta de menor resistência de saída é mais rígida.}
\end{questao}

\begin{questao}[4]{}
Uma caracterização experimental produz pontos que não ficam sobre uma reta. Dê três causas físicas e indique como separar erro de medição de não linearidade real.
\resposta{Aquecimento, limitação de corrente/saturação e elementos não lineares; repetir em níveis menores e com mais pontos}
\resolucao{Erros aleatórios geram resíduos sem padrão; não linearidade tende a produzir curvatura sistemática. Repetir o ensaio com amplitudes menores ajuda a revelar aquecimento, saturação ou mudança de estado.}
\end{questao}

\begin{questao}[4]{}
Uma rede será usada com vinte cargas diferentes. Compare a estratégia de resolver o circuito completo a cada carga com a de obter uma vez $V_{Th}$ e $R_{Th}$.
\resposta{O equivalente exige uma caracterização inicial e depois apenas uma expressão escalar por carga; a análise completa repete o sistema vinte vezes}
\resolucao{O ganho cresce com o número de cargas. Thévenin é especialmente vantajoso quando a rede interna permanece fixa e apenas a carga varia; perde vantagem se forem necessárias grandezas internas da rede.}
\end{questao}

